\section{Related Work}
\label{sec:related}
\textbf{Heterogeneous graph learning.} Type-aware message passing was
popularized by attention-based encoders that learn a distinct transformation
per node/edge type~\cite{tesfaye2019han}, later scaled to
billion-edge industrial graphs via type-partitioned
sampling~\cite{oyelaran2023heterogeneous}. These methods assume a fixed
snapshot and do not model recency.

\textbf{Temporal graph learning.} Sequence- and attention-based temporal
encoders~\cite{brandt2022temporal,marceau2021dysat} model how a node's
representation drifts as new edges arrive, typically through a memory
module or a sliding attention window over recent interactions. A survey by
Nakashima and Ferreira~\cite{nakashima2021survey} catalogs over forty
variants but notes that fewer than a fifth handle more than one node type.

\textbf{Contrastive and self-supervised objectives.} Contrastive
pretraining, first popularized for images~\cite{yamazaki2021simclr}, was
adapted to graphs by treating augmented views of the same subgraph as
positive pairs~\cite{alavi2020contrastive}. \sysname{} departs from
augmentation-based positives, instead using temporally adjacent real
interactions, which avoids hand-tuned augmentation policies.

\textbf{Scalability and evaluation.} Petrosyan and
Deveaux~\cite{petrosyan2022scalable} study sampling strategies for
billion-edge training, and Husejnovic and
Delacroix~\cite{husejnovic2020fraud} apply graph encoders to industrial
fraud detection, both motivating our efficiency-focused design. Osei and
Lindgren~\cite{osei2022benchmark} highlight reproducibility gaps in
temporal-graph benchmarks, which informed our choice to release full
hyperparameters in Section~\ref{sec:experiments}.
